\chapter{Introduction}
\section{Background}
RedSentinel is an AI-augmented web application security scanner purpose-built to detect Cross-Site Scripting (XSS) vulnerabilities. XSS remains one of the most persistently exploited vulnerability classes in modern web applications. According to the OWASP Top 10, injection-related vulnerabilities have appeared consistently across multiple editions of the industry-standard web application risk rankings, with Cross-Site Scripting discussed within A03:2021 (Injection) as a prominent client-side security concern \cite{owasp_top10_2021}. The widespread prevalence of XSS is further corroborated by the Common Vulnerabilities and Exposures (CVE) database, where XSS routinely accounts for a substantial share of annually reported web vulnerabilities \cite{mitre_cve}.

Traditional XSS scanners operate on static, manually curated payload lists paired with simple string-matching heuristics. While adequate for detecting trivial reflections, these approaches fail in two important respects: they are blind to injection context meaning they cannot distinguish between a parameter reflected inside a JavaScript string literal and one reflected inside an HTML attribute value, contexts that demand entirely different payload structures and they are easily defeated by modern Web Application Firewalls (WAFs), which signature-match against known payload patterns \cite{bau2010state}.

RedSentinel addresses both limitations through a five-phase pipeline that integrates machine learning-based injection context classification at the core of its analysis workflow. The context classification is performed by a fine-tuned DistilBERT model, a distilled variant of the BERT transformer architecture, introduced by Devlin et al.\cite{devlin2019bert} and subsequently compressed by Sanh et al.\cite{sanh2019distilbert} to approximately 40\% of BERT's parameter count while retaining 97\% of its language understanding capability on the GLUE benchmark. The model is fine-tuned on a domain-specific dataset of 59,122 labeled XSS payloads, enabling it to classify injection contexts (HTML body, HTML attribute, JavaScript string, JavaScript code block, URL context, CSS context, and DOM sink) from source code fragments surrounding the injection point.

The system's architecture follows a polyglot microservices pattern. A NestJS/TypeScript orchestration core \cite{nestjs} manages the HTTP API, WebSocket real-time event streaming via Socket.io \cite{socketio}, job queue scheduling via BullMQ backed by Redis \cite{bullmq}, and PostgreSQL persistence via TypeORM \cite{typeorm}. Three independently deployable Python FastAPI microservices \cite{fastapi} handle the AI-intensive concerns: context analysis (port 5001), payload generation with mutation and obfuscation (port 5002), and HTTP/headless-browser fuzzing via Playwright \cite{playwright} (port 5003). Asynchronous job execution through BullMQ ensures the REST API remains non-blocking while scans execute in background workers, and Socket.io WebSocket events deliver real-time progress and vulnerability findings to the Next.js dashboard \cite{nextjs}.

The payload generation subsystem draws on a curated bank of 59,122 labeled XSS payloads. Payload selection is context-aware and WAF-evasion-weighted, followed by a mutation engine that applies systematic transformations (case variation, Unicode and HTML entity encoding, whitespace manipulation, and event-handler alternation) and an obfuscation layer that applies higher-order JavaScript transforms such as Base64/atob encoding, String.fromCharCode construction, and eval indirection. This approach is consistent with established research on WAF bypass techniques, which demonstrates that encoding diversity and structural variation significantly reduce detection rates of signature-based filters \cite{bau2010state}.

Vulnerability confirmation in the fuzzing phase employs a two-stage verification strategy: HTTP response reflection analysis followed by headless Chromium execution via Playwright, using a canary window property injection technique to confirm actual JavaScript execution rather than mere textual reflection. This dual-verification approach is aligned with best practices in automated XSS confirmation research, which distinguishes passive reflection (a necessary but insufficient condition for XSS) from confirmed execution. The system concludes each scan by generating vulnerability reports in HTML, PDF (via Puppeteer \cite{puppeteer}), and JSON formats, suitable for both human review and CI/CD pipeline integration.

\section{ Problem Statement}
Although XSS scanners and WAFs can identify many basic vulnerabilities, they often depend on static payloads and basic reflection checks, making them less effective at handling different injection contexts, WAF/filter bypasses, browser execution verification, and structured reporting.

\section{ Objectives}
The primary objectives of the RedSentinel are:
\begin{enumerate}[itemsep=2pt, topsep=-15pt]
    \item To design and implement RedSentinel, an AI-assisted, context-aware XSS scanner that crawls targets, analyses injection contexts, generates and ranks payloads, fuzzes and verifies vulnerabilities.
    \item To generate structured, actionable vulnerability reports in multiple formats.
\end{enumerate}
